\documentclass[tcc,capa,table]{textsi}

%\usepackage[brazilian]{babel} % Linguagem do documento: Português-brasileiro
\usepackage[utf8]{inputenc} % Codificação do arquivo .tex. IMPORTANTE: O arquivo deve ser salvo utilizando a codificação UTF-8 para mostrar corretamente os caracteres!
\usepackage[T1]{fontenc} % Permite que se faça cópia do texto do PDF de maneira correta, usando a codificação de fonte T1

% Alguns pacotes disponíveis no latex para adição e manipulação de elementos do documento
\usepackage{array,booktabs}
\usepackage{enumitem}
\usepackage{listings}
\usepackage{makecell}
\usepackage{graphicx,subfigure}
\usepackage{multirow}
\usepackage{xcolor}
\usepackage{pdfpages}
\usepackage{lipsum}
\usepackage{longtable}
\usepackage{placeins}
\usepackage{float}
\usepackage{xcolor}
\graphicspath{{images/}}

\lstset{
  basicstyle=\ttfamily\footnotesize,
  breaklines=true,
  frame=single,
  columns=fullflexible,
}

% Título do trabalho
\title{nexo - Um aplicativo para rastrear flutuações de humor para pessoas com TAB}

% Autor do trabalho. Primeiro argumento: Último sobrenome; Segundo argumento: Restante do nome
\author{Porcellis}{Pedro Lucas}

% Orientador do trabalho. Primeiro argumento (opcional): titulação do orientador. Segundo argumento (obrigatório): Último sobrenome; Terceiro argumento (obrigatório): Restante do nome
\advisor[Prof.~Me.]{Nunes}{Flávio}

% Instituição do orientador do trabalho.
\instituicaoadvisor{IFSul}

% Descomente a linha abaixo se tens co-orientador. Segue o mesmo formato do orientador do trabalho.
%\coadvisor[Prof.~Dr.]{Aguiar}{Marilton Sanchotene de}

% Descomente a linha abaixo se tens colaborador. Segue o mesmo formato do orientador do trabalho.
%\collaborator[Prof.~Dr.]{Aguiar}{Marilton Sanchotene de}

% Membro 1 da banca. Um argumento único, contendo titulação e nome completo.
\membroi{Prof.~Dr. Fulano de Souza e Silva}
% Instituição do membro 1 da banca.
\instituicaomembroi{UFI}

% Membro 2 da banca. Segue o mesmo formato que o membro 1.
\membroii{Prof. Dr. Sicrano Silveira Rosa}
% Instituição do membro 2 da banca.
\instituicaomembroii{UFPG}

% Descomente se tens um terceiro membro na banca. Membro 3 da banca. Segue o mesmo formato que o membro 1.
\membroiii{Prof. Dr. Beltrano Pereira Tavares}
% Descomente se tens um terceiro membro na banca. Instituição do membro 3 da banca.
\instituicaomembroiii{UFCR}

% Palavras-chave do trabalho desenvolvido.
\keyword{Desenvolvimento de Aplicativo}
\keyword{Transtorno Afetivo Bipolar}
\keyword{React Native}

% XXX: Fix longtables not having a prefix. Fucking POS template; shoulda stick
% to rawdogging ABNT2 or just use the IFRJ one which is actually complete

\makeatletter
\def\LT@c@ption#1[#2]#3{%
  \LT@makecaption#1\fnum@table{#3}%
  \def\@tempa{#2}%
  \ifx\@tempa\@empty\else
    \addcontentsline{lot}{table}%
      {\protect\numberline{\fnum@table}{\ignorespaces #2}}%
  \fi}
\makeatother

\begin{document}

% Descomente se você tem orientadora, e não orientador.
%\renewcommand{\advisorname}{\orientadora}

% Descomente se você tem coorientadora, e não orientador.
%\renewcommand{\coadvisorname}{\coorientadora}

% Faz com que o título apareça
\maketitle

\sloppy

% Ficha catalográfica e folha de aprovação
\fichacatalografica
\folhadeaprovacao

%Opcional - dedicatória
\begin{dedicatoria}
	Texto da Dedicatória
\end{dedicatoria}

%Opcional - agradecimentos
\begin{agradecimentos}
	Texto dos Agradecimentos
\end{agradecimentos}

%Opcional - epígrafe
\begin{epigrafe}
	"Texto da Epígrafe" \\
	Autor da Epígrafe
\end{epigrafe}

% Resumo em língua vernácula (a que você está escrevendo o documento)
% TODO: Rework this a bit more
% TODO: 150 words min, not more than one page, obviously, always saying that it
% concluded the proposed objectives
\begin{abstract}
  Esse projeto aborda o desenvolvimento de um aplicativo de auxílio e controle
  para pessoas com transtorno de afetivo bipolar (TAB). A proposta é uma
  aplicação cujo usuário pode utilizar para monitorar o registro diário tanto
  de mudanças no humor quanto também no controle e registro de intervenções,
  qualidade de sono e gatilhos, com o objetivo de facilitar a identificação da
  fase cíclica da TAB. O trabalho abrange desde a construção da pesquisa
  inicial, o desenvolvimento da aplicação e testes. O sistema é praticamente
  inteiramente desenvolvido na linguagem de programação JavaScript, tanto com o
  backend/API e a aplicação, desenvolvida em React Native.
\end{abstract}

% Resumo em língua estrangeira
\begin{englishabstract}%
 	{Mobile application, Bipolar Affective Disorder, Mood Tracker}

  This work addresses the development of a mood tracking mobile application
  aimed at people with Bipolar Affective Disorder (BAP). The proposal is a
  daily mood tracker, which also has modules for tracking sleep quality, which
  can aid in identifying hypomanic or manic cycles. It’ll also design a
  tracking of other elements such as interventions, sleep patterns, triggers,
  and reminders, which engages in a more faithful detection of cycles and
  patterns. The work itself goes all the way from the initial research up to
  the development of the app, layout, and tests. The system itself is mostly
  written in JavaScript.

\end{englishabstract}

%Lista de Figuras
\listoffigures

%Lista de Tabelas
\listoftables

% Lista de códigos
\lstlistoflistings

%lista de abreviaturas e siglas
\begin{listofabbrv}{REST} % <- Insira aqui a maior sigla da sua lista de siglas.
	% O conteúdo dentro dos colchetes é a sigla, e o texto ao lado é o seu significado
	\item[TAB] Transtorno Afetivo Bipolar
  \item[BAP] Bipolar Affective Disorder
  \item[TCC] Terapia Cognitivo-comportamental
  \item[UML] \textit{Unified Modeling Language}
  \item[DAO] \textit{Data Access Object}
  \item[REST] \textit{Representational State Transfer}
  \item[API] \textit{Application Programming Interface}
  \item[UX] \textit{User Experience}
  \item[SSH] \textit{Secure Shell}
  \item[CI] \textit{Continuous Integration}
  \item[CD] \textit{Continuous Deployment}
\end{listofabbrv}

%Sumario
\tableofcontents

% Chapter é a seção primária do documento. Toda vez que é chamado este comando, uma quebra de página é adicionada;
\input{chapters/introduction}
\chapter{Saúde e autocuidado mental no mundo digital}

% TODO: Move these refs into bibtex

No quesito de saúde mental, a bipolaridade ou pelo seu nome científico,
Transtorno Afetivo Bipolar (TAB), é uma das doenças mentais mais comuns e
responsável por uma gama muito grande de disfunções mentais. De acordo com
o estudo conduzido por \citet{Merikangas2011}, no âmbito do \textit{World Mental
Health Survey Initiative} da Organização Mundial da Saúde, com mais de 61 mil
adultos entrevistados em 11 países, a prevalência ao longo da vida do espectro
bipolar é de 2,4\% da população mundial, e a prevalência em um período de 12
meses é de 1,5\%. No contexto do Brasil, segundo a Associação Brasileira de
Transtorno Bipolar (ABTB), capítulo brasileiro da \textit{International Society for
Bipolar Disorders} (ISBD) \citep{abtb}, a bipolaridade afeta cerca de
4\% da população, o que representaria uns 7,6 milhões de
pessoas. Esse número, embora soe pouco expressivo, levando em consideração o
contingente de pessoas no país, é preocupante, dado o fato de que a
bipolaridade é (1) complexa de se detectar e (2) muitas vezes, acaba passando
despercebida ou enquadrada dentro de um quadro de depressão unipolar
(depressão).

Embora o autorregistro seja amplamente recomendado no acompanhamento
de pessoas com TAB, dado seu potencial de validade clínica e de efeito
sobre desfechos relevantes~\cite{Faurholt:2016}, sua realização de
maneira contínua representa um desafio.
Em muitos casos, o paciente depende da própria memória para reconstruir
retrospectivamente seu estado emocional, a qualidade do sono,
acontecimentos cotidianos e outros fatores que podem ter influenciado
seu humor. Esse processo está sujeito, naturalmente, a vieses de
recordação e dificulta a identificação de padrões temporais, especialmente
quando as informações são coletadas apenas durante consultas clínicas.

Nesse contexto, o \textit{Ecological Momentary Assessment} (EMA)
consolidou-se como uma importante abordagem metodológica para coleta
de dados em saúde mental. Diferentemente de avaliações retrospectivas,
o EMA propõe que informações sejam registradas no ambiente natural do
indivíduo e próximas ao momento em que as experiências ocorrem, reduzindo
vieses de memória e permitindo observar a dinâmica dos estados afetivos
ao longo do tempo \citep{EbnerPriemer2009}. Além de proporcionar maior
validade ecológica aos dados coletados, essa estratégia tem sido
empregada em diversas áreas da pesquisa comportamental, incluindo
transtornos do humor \citep{Collins2003, EbnerPriemer2009}, permitindo
compreender a relação entre fatores contextuais, eventos cotidianos e
alterações do humor.

Nesse contexto, o uso de tecnologias digitais tem se mostrado uma via
promissora para sistematizar esse tipo de acompanhamento. Pesquisas indicam
que pessoas com transtornos do espectro bipolar demonstram alta adesão ao
monitoramento emocional quando este é feito de forma acessível, visual e
integrada ao seu cotidiano, com estudos reportando taxas de conclusão acima
de 90\% em protocolos de avaliação momentânea
via \textit{smartphone} \citep{Ryan2020} \footnote{
  Vale citar porém, que esse artigo tem limitações de generalização:
  os resultados são limitados a um grupo de pessoas que possuiam um maior
  contato com tecnologia.
}.

Entre as formas mais comuns de implementação do EMA encontra-se o
autorregistro contínuo (\textit{journaling}), no qual o indivíduo
registra sistematicamente aspectos relevantes do cotidiano, como
humor, qualidade do sono, eventos significativos, sintomas percebidos
e outras informações contextuais. Com o acúmulo desses registros,
torna-se possível identificar padrões de comportamento, reconhecer fatores
associados às alterações de humor e favorecer o desenvolvimento da
autoconsciência sobre a própria condição, constituindo uma ferramenta
complementar ao processo de autoconhecimento e autogestão do
comportamento \citep{Melton2013}.

A ampla disseminação dos dispositivos móveis ampliou significativamente
as possibilidades de utilização dessas estratégias fora do ambiente
clínico. Smartphones acompanham o indivíduo durante praticamente toda
sua rotina, permitindo que registros sejam realizados no momento em que
as experiências ocorrem e reduzindo o esforço necessário para manutenção
do acompanhamento contínuo. Como consequência, diversas aplicações
voltadas ao monitoramento da saúde mental passaram a incorporar princípios
do EMA em seus fluxos de utilização, oferecendo registros estruturados
e visualizações históricas capazes de auxiliar usuários e profissionais
de saúde \citep{Faurholt:2016}.

Apesar da crescente disponibilidade de aplicações voltadas ao acompanhamento
do humor, estudos recentes que avaliaram especificamente os aplicativos mais
utilizados por pessoas com TAB para monitoramento de humor e sono apontam
limitações importantes quanto à integração entre esses domínios, ao controle
do usuário sobre seus próprios dados e à acessibilidade das políticas de
privacidade, muitas vezes redigidas em nível de leitura incompatível com o
público geral \citep{Morton2022}.

Em paralelo, análises mais amplas do ecossistema de aplicativos de
saúde mental também identificam riscos recorrentes de privacidade,
como permissões desnecessárias e vazamento de dados sensíveis \citep{Iwaya:22}.
Essas limitações evidenciam a necessidade de soluções que conciliem
boas práticas de engenharia de software, proteção de dados pessoais e
recursos que apoiem a identificação de padrões emocionais ao longo do tempo.

Diante desse contexto, este trabalho propõe o desenvolvimento de uma
aplicação móvel fundamentada nos princípios do
\textit{Ecological Momentary Assessment} voltada ao acompanhamento
cotidiano de pessoas com Transtorno Afetivo Bipolar. O sistema busca
integrar registros estruturados de humor, qualidade do sono, gatilhos
e ações de cuidado em uma única plataforma, produzindo indicadores que
auxiliem o usuário na visualização de padrões e no desenvolvimento da
consciência sobre seu próprio estado afetivo. Ressalta-se que a aplicação
não possui finalidade diagnóstica nem substitui o acompanhamento realizado
por profissionais de saúde, sendo concebida como uma ferramenta de apoio
ao autorregistro e ao autocuidado.

\chapter{Metodologias}

Como primeiro passo, para conseguir viabilizar construir um \textit{software}
com qualidade, é necessário adotar metodologias que fundamentam a ideia e o
processo. De tal sorte que, ajudem a levantar requisitos, priorizar
funcionalidades, quebrar o escopo em entregáveis que possam ser validados, e
desenvolver o aplicativo/sistema de forma sustentável. Esse capítulo se
debruça, tanto na abordagem da metodologia de pesquisa do projeto, quanto
também detalha o método de desenvolvimento de \textit{software} implementado.

\section{Metodologia de Desenvolvimento de \textit{software}}

Existem diversos métodos de desenvolvimento, cada um com seus respectivos prós
e contras. Todos tentam, de sua forma, responder à necessidade de como dar
forma ao processo de desenvolvimento, para que ele flua e assegure que os
requisitos levantados sejam atendidos, que os membros envolvidos no processo,
tanto de desenvolvimento em si, quanto dos atores que o demandam tenham seus
anseios respondidos de forma minimamente transparente e satisfatória.

Além disso, a escolha de uma boa metodologia, adequada para as necessidades do
sistema, permite que os prazos de entrega sejam melhor controlados e mais
precisos \cite{pressman:2021}, minimizando a ocorrência de falhas e gestão do
tempo.

As ideias fundamentais que norteiam esse projeto são baseadas em dois conceitos
chaves da engenharia de \textit{software}: desenvolvimento incremental e
desenvolvimento de \textit{software} ágil. O IEEE define que a engenharia de
\textit{software} é:

\begin{quote}
  Engenharia de \textit{software}: A aplicação de uma abordagem sistemática,
  disciplinada e quantificável no desenvolvimento, na operação e na manutenção de
  \textit{software}; isto é, a aplicação de engenharia ao \textit{software}.
  \cite[p.~67, tradução nossa]{ieee:1990}
\end{quote}

Portanto, a engenharia de \textit{software} propõe estruturar a criação de
\textit{software} de forma mais organizada e eficiente. Pressman vai dar um
destaque para os modelos de desenvolvimento de \textit{software} sequenciais e
incrementais que moldam boa parte das práticas de engenharia de
\textit{software} nas últimas décadas.

O modelo em cascata (\textit{waterfall}) é um modelo sequencial, dos mais
tradicionais e antigos a existirem, cujo qual necessita que cada etapa --
levantamento de requisitos, projeto, implementação, testes -- estejam
concluídas antes da próxima etapa começar. Embora forneça uma estrutura rígida
e lógica, é frequentemente criticado pela baixa adaptabilidade em relação às
mudanças durante o processo de desenvolvimento.

Em outro âmbito, o modelo incremental, trabalha com a premissa de que um
desenvolvimento, feito em pequenas porções, agrupam-se para totalizar o projeto
em si. Trabalho a concepção de incrementos funcionais, permitindo uma maior
flexibilidade com as mudanças no processo ao todo. Também ajuda a iterar
rápidamente em cima do desenvolvimento do sistemas, criando um produto mínimo
viável, que possa preencher a proposta, usando o mínimo de tempo, reduzindo as
funcionalidades ao seu estado basal. Essa metodologia também ajuda a reduzir
custos e riscos, já que necessita que as funcionalidades estejam em seu estado
mínimo, porém funcionais.

Outros modelos tradicionais citados por Pressman incluem o modelo de
prototipação e o modelo em espiral, o primeiro trabalhando na ideia de
elaboração de protótipos pequenos que validem a hipótese, e depois passe para a
implementação final; e o modelo em espiral une a natureza iterativa da
prototipação com os aspectos sistemáticos do modelo em cascata, sendo que acaba
por ser um modelo geralmente usado para desenvolvimento de aplicações em larga
escala, visto que oferece mais segurança ao processo de desenvolvimento. Apesar
da influência histórica, tais modelos acabam por serem um tanto quanto
burocráticos e em um projeto com um escopo de equipe reduzida e onde necessita
um nível de flexibilidade maior, como este, portanto, uma adoção de uma
metodologia mais flexível se faz necessária.

Pressman, comenta que, por meados de 2001, um grupo de desenvolvedores
deflagrou o "Manifesto Ágil", um documento desenhando uma nova metodologia de
desenvolvimento, destacando a satisfação dos patrocinadores do sistema,
entregas mais rápidas, flexibilidade durante o desenvolvimento do
\textit{software} de forma a evitar que os desenvolvedores ficassem empacados e
reduzissem rigidez. Os principais destaques da abordagem ágil é a proposta de
que a mesma não é escrita em pedra, ou seja, permite flexibilidade no como
aplicá-la. Nesse cenário, existem algumas abordagens sobre ela, como o
\textit{Kanban} e o \textit{Scrum}, que dão um pouco mais de forma a ideia
apresentada com o desenvolvimento ágil o que casa relativamente bem com o
modelo de desenvolvimento incremental.

Com base nesse panorama, o presente trabalho adota uma abordagem baseada na
filosofia ágil, estruturada por meio da utilização de um quadro de Kanban,
adaptado à realidade de um desenvolvimento individual. Diferentemente de
modelos tradicionais que envolvem um maior formalismo, o \textit{Kanban}
permite visualizar e gerenciar o fluxo de trabalho de forma contínua, com foco
na entrega incremental de funcionalidades. Por ser um método não prescritivo, o
\textit{Kanban} oferece liberdade para que as tarefas possam ser priorizadas,
ajustadas ou redefinidas conforme a evolução do desenvolvimento, característica
especialmente útil quando há apenas um desenvolvedor responsável por todas as
etapas do processo. Vale ressaltar também, que a metodologia adotada acaba por
ajudar a manter o fluxo de desenvolvimento contínuo sem sobrecarga, já que
limita a quantidade de tarefas em progresso, gerando um movimento de maior
produtividade e entregas.

A escolha pelo \textit{Kanban}, portanto, oferece um equilíbrio sustentável
entre estrutura e flexibilidade, permitindo que o desenvolvimento avance de
maneira organizada, e sem os entraves burocráticos dos modelos mais rígidos.
Essa abordagem permite que o projeto se beneficie dos princípios ágeis, tais
como a entrega contínua de valor, adaptação a mudanças e foco na simplicidade,
ao mesmo tempo que mantém o controle sobre o progresso e a qualidade do produto
final.

\section{Metodologia de Pesquisa}

Entende-se por pesquisa científica, a investigação de um fenômeno, de forma a
encontrar a solução para algum problema \cite{Coelho:2020}. Portanto, a
metodologia científica se torna o conjunto de procedimentos para esta
investigação, cujo qual perpassa a coleta e análise de dados. Neste capítulo,
abordaremos as metodologias adotadas e suas respectivas justificativas, para
além dos dados encontrados para tal.

\subsection{Tipo de Pesquisa}

Segundo \cite{Lakatos:2022}, a pesquisa científica pode ser classificada
quanto aos seus objetivos como exploratória, descritiva ou explicativa. Este
trabalho caracteriza-se principalmente como uma pesquisa aplicada de natureza
exploratória, uma vez que busca investigar práticas, limitações e necessidades
relacionadas a ferramentas digitais para acompanhamento emocional, com o
objetivo de desenvolver uma solução tecnológica prática.

Conforme Jung (2009), a pesquisa aplicada visa gerar conhecimento para
aplicação direta em problemas específicos, com interesse prático e utilitário,
sendo comum em projetos de desenvolvimento de \textit{software}, especialmente
nas áreas de Sistemas de Informação e Ciência da Computação, nas quais o
produto final visa resolver uma demanda concreta do público-alvo.

\subsection{Instrumento de coleta de dados}

Como instrumento de coleta de dados, optou-se pela aplicação de um questionário
estruturado, elaborado com o intuito de compreender hábitos, preferências e
resistências de usuários em relação ao uso de aplicativos voltados para o
acompanhamento emocional e práticas de autorreflexão.

A escolha do questionário se justifica pela sua abrangência, agilidade na
coleta de dados e facilidade de análise quantitativa e qualitativa. Além disso,
como a amostra é composta por usuários potenciais, a aplicação do questionário
de forma digital (via um formulário online) permite alcançar um público mais
amplo com baixo custo e maior comodidade, como também reforça Jung (2009).

O questionário, disponível no Apêndice B, foi bem recebido, e obteve por volta
de 52 respostas válidas. Como a divulgação foi de forma aberta, existia uma
pergunta inicial com uma função chaveadora, a fim de verificar se os
respondentes se encaixavam no perfil de potenciais usuários. Isso se justifica
pela natureza sensível do tema, que envolve saúde mental e experiências
subjetivas.

Na pergunta ``Você já teve algum interesse ou experiência com o monitoramento do
seu humor ou estado emocional?'', apenas 8\% afirmaram que não tem e não teriam
interesse em realizar esse tipo de acompanhamento. Esse dado corrobora outro
resultado relevante: 53\% dos participantes afirmaram já ter, em algum momento
da vida, realizado algum tipo de monitoramento de humor.

Em outros horizontes, mais da metade dos participantes expressou interesse em
acompanhar como o humor varia ao longo do tempo, muitas vezes com sugestões
explícitas como: “Seria legal ver um gráfico semanal ou mensal do meu humor”.
Tal observação reforça a receptividade dos participantes em relação à ideia de
detecção de padrões emocionais, sugerindo que esse tipo de funcionalidade
representa um diferencial positivo frente aos aplicativos existentes.

É notável que os participantes demonstram uma preocupação recorrente com a
privacidade dos dados coletados. Aproximadamente 67\% dos participantes
indicaram desconforto com a exposição de suas informações em plataformas que
não explicam de forma clara como os dados são usados ou armazenados. Essa
percepção está alinhada com a literatura sobre saúde digital. Por exemplo,
\cite{Iwaya:22}, argumentam que a ausência de mecanismos transparentes de
consentimento, controle e anonimização pode minar a confiança dos usuários em
aplicativos de bem-estar e monitoramento contínuo de saúde, especialmente em
contextos sensíveis como saúde mental.

Outro ponto importante revelado pelo levantamento é o desejo por um tom mais
humano, subjetivo e acolhedor. Cerca de 52\% dos participantes relataram
incômodo com ferramentas que adotam uma linguagem genérica, motivacional ou
automatizada demais. Isso reforça a necessidade de uma interface mais empática,
que valorize tanto a escrita livre quanto a análise assistida. Também foram
mencionadas sugestões de funcionalidades consideradas excepcionais, como
inferência automática de ciclos de humor (66\% demonstraram interesse),
personalização da forma de visualização das informações (como gráficos ou
comparativos), e integração com dados contextuais ou ambientais, reforçando a
ideia de um sistema que vá além do simples diário, oferecendo inteligência
emocional útil, mas sem perder a sensibilidade que o tema exige.

\subsection{Análise de Similares}

% TODO: reescrever, arrumar citação
A metodologia trata-se de uma avaliação sobre os sistemas relacionados ao tema
proposto, observando e levantando desafios e possibilidades, validação de
\textit{softwares} e a satisfação dos usuários (Martins, 2018). Seguindo assim,
foram pesquisados sistemas similares encontrados na internet com o objetivo de
realizar uma revisão crítica-reflexiva, analisando os aplicativos que oferecem
algum tipo de diário digital de humor, ou permite algum acompanhamento
reflexivo de variações de emoções.

Para realizar esse processo, de forma a simular também um usuário, foi
pesquisado na Google Play Store \footnote{Google Play Store — Disponível em:
https://play.google.com}, por \textit{mood tracker}, retornando centenas de
resultados. Nesse quesito, estabeleceu os seguintes critérios objetivos: (1)
pode ser um aplicativo em português ou inglês; (2) o aplicativo tem que
apresentar algum tipo de gráfico das entradas feitas (3) tem que possuir
entrada textual.

\subsubsection{Baseline}

O Baseline é um aplicativo de código aberto voltado ao acompanhamento do
bem-estar emocional por meio de uma abordagem minimalista. A ferramenta permite
o registro diário de humor utilizando entradas textuais ou gravações de áudio,
além de disponibilizar notificações personalizadas e revisões periódicas dos
registros realizados.

\begin{figure}[H]
\centering
\caption{Aplicativo Baseline}
\label{fig:screenshot-baseline}
\includegraphics[width=0.40\textwidth,height=0.40\textheight,keepaspectratio]{screenshot-baseline.jpeg}
\source{O autor}
\end{figure}

Entre os aspectos observados, destaca-se a natureza open-source da aplicação,
característica que favorece a transparência e a possibilidade de adaptação por
terceiros. Sua interface apresenta baixo nível de complexidade e reduz a
quantidade de interações necessárias para o registro das informações.
Entretanto, nota-se que o sistema não oferece visualizações temporais avançadas
nem mecanismos de personalização mais aprofundados. Além disso, o
acompanhamento emocional é realizado principalmente por meio de uma pontuação
de humor, sem identificar ciclos emocionais ou relacionar registros a gatilhos
e intervenções específicas.

\subsubsection{DailyBeans}

O DailyBeans é um dos aplicativos de monitoramento de humor mais populares
disponíveis atualmente, possuindo como principal característica a simplicidade
de uso. Nota-se que  possui uma linguagem bem lúdica e informal, por exemplo, o
sistema utiliza o conceito de ``feijões'' representando o estado de humor do
usuário. Destaca-se que o aplicativo permite selecionar categorias de
sentimentos e adicionar atividades para preencher melhor o registro do estado
emocional do usuário para além de um "Feliz", "Triste", etc. Existe também uma
bastante vasta opções de categorias, cada qual com uma ilustração que
representa o que aquilo representa. O aplicativo também apresenta "Temas
sazonais", permitindo que o usuário customize cores, icones, e demais
informações. Um elemento interssante é que as informações adicionadas também
podem ser \textit{offlines}, até uma certa quantidade de entradas/registros,
necessitando de uma conta para poder continuar salvando os dados informados.

\begin{figure}[H]
\centering
\caption{Aplicativo DailyBeans}
\label{fig:screenshot-dailybeans}
\includegraphics[width=0.40\textwidth,height=0.40\textheight,keepaspectratio]{screenshot-dailybeans.jpeg}
\source{O autor}
\end{figure}

Entre suas funcionalidades, além do registro de humor e categorização de
atividades e sentimentos, o aplicativo conta com gráficos e históricos. Uma
interface bastante coesa e visualmente agradável. Um outro elemento é que o
aplicativo permite relacionar indicadores de sono com bastante precisão, como
poder indicar o momento que foi dormir, momento que acordou, anotações, etc. O
aplicativo permite adicionar pelo menos uma foto no plano gratuito, além de
indicar elementos como ``Qual o clima do dia'', se o usuário esteve com
``amigos'', ``familiares''.

Observa-se que o aplicativo possui código fechado e restringe parte de suas
funcionalidades mais avançadas à versão paga ou compras na Loja interna do
aplicativo para obter por exemplo, personalização de icones, adicionar mais
fotos. Além disso, o foco em registros rápidos reduz a flexibilidade para
inserção de reflexões textuais mais detalhadas, e uma exibição excessante de
anúncios atrapalha a utilização.

\subsubsection{DailyMood}

DailyMood é um aplicativo voltado ao registro cotidiano do estado emocional,
oferecendo uma solução enxuta para acompanhamento de humor. Sua proposta
baseia-se em uma interação mínima, permitindo que o usuário selecione
rapidamente o sentimento predominante do dia, sem a necessidade de descrições
extensas. Os registros são organizados em um calendário, o que facilita a
leitura histórica e a identificação de padrões emocionais ao longo do tempo.

\begin{figure}[H]
\centering
\caption{Aplicativo DailyMood}
\label{fig:screenshot-dailymood}
\includegraphics[width=0.40\textwidth,height=0.55\textheight,keepaspectratio]{screenshot-dailymood.jpeg}
\source{O autor}
\end{figure}

Apesar dos recursos disponíveis, observa-se que é uma solução bastante simples.
Não existe outros recursos além de informar o humor predominante e uma
anotação. Um fator interessante é que não necessita de contas, rodando
completamente offline no dispositivo do usuário. Adicionalmente, apresenta
limitações relacionadas à exportação dos dados registrados e à personalização
das informações armazenadas.

\subsubsection{Reflectly}

O Reflectly é um aplicativo de diário pessoal que utiliza recursos de
inteligência artificial para promover práticas de autorreflexão e
acompanhamento emocional. Sua proposta enfatiza a experiência do usuário e a
construção de narrativas pessoais a partir dos registros realizados.

Entre suas funcionalidades estão entradas textuais livres, perguntas guiadas
diárias, análises emocionais automatizadas e personalização visual da
interface. O sistema apresenta uma experiência moderna e intuitiva,
incentivando o uso contínuo por meio de elementos de engajamento.
Um ponto interessante é o fato do aplicativo separar o ``humor'' dos
``sentimentos'', permitindo que os usuários avaliem ambos de forma
independente.

Como pontos observados, destaca-se o fato de ser uma solução proprietária e
dependente de coleta significativa de dados para alimentar seus mecanismos de
análise. Outro ponto negativo é o fato de que o aplicativo apresenta
elementos como estatísticas e histórico, porém necessita de uma conta
paga para acessar esses recursos.

Além disso, sua abordagem está mais voltada ao diário pessoal do que
ao acompanhamento sistemático de indicadores de saúde mental.

\begin{figure}[H]
\centering
\caption{Aplicativo Reflectly}
\label{fig:screenshot-reflectly}
\includegraphics[width=0.40\textwidth,height=0.55\textheight,keepaspectratio]{screenshot-reflectly.jpeg}
\source{O autor}
\end{figure}

\subsubsection{eMoods}

O \textit{eMoods} é um aplicativo de registro de humor especificamente para
pessoas com TAB. O aplicativo tem uma interface bem compacta, bastante
poluída visualmente, em estilo formulário de entrada de dados.

O aplicativo consegue registrar dados de sono, além de perguntas como
ansiedade, estado depressivo, e outras coisas como ``você praticou higiene
pessoal?'', ``você tomou medicamentos?'', ``você se alimentou bem?'' e afins.
Um ponto positivo é que o aplicativo também permite exportar os dados em
formato CSV e gerar relatórios em formato PDF.

Entretanto, a aplicação apresenta limitações relacionadas à experiência do
usuário, possuindo interface considerada desatualizada quando comparada a
soluções contemporâneas. Também não oferece suporte para registros textuais
extensos ou entradas por voz, além de ser muito confuso de usar.

\begin{figure}[H]
\centering
\caption{Aplicativo eMoods}
\label{fig:screenshot-emoods}
\includegraphics[width=0.40\textwidth,height=0.55\textheight,keepaspectratio]{screenshot-emoods.jpeg}
\source{O autor}
\end{figure}

\FloatBarrier

\subsection{Análise comparativa}

\begin{table}[H]
    \centering
    \caption{Comparativo de funcionalidades entre aplicativos analisados}
    \label{tab:comparativo_apps}
    \begin{tabular}{lccccc}
        \hline
        Funcionalidade &
        Baseline &
        DailyBeans &
        DailyMood &
        Reflectly &
        eMoods \\

        \hline
        Registro de humor & Sim & Sim & Sim & Sim & Sim \\
        Registro de sono & Não & Não & Não & Não & Sim \\
        Texto livre & Sim & Parcial & Sim & Sim & Sim \\
        Gráficos & Não & Sim & Sim & Sim & Sim \\
        Gatilhos/intervenções & Não & Parcial & Não & Não & Parcial \\
        Exportação de dados & Sim & Parcial & Parcial & Parcial & Sim \\
        Código aberto & Sim & Não & Não & Não & Não \\
        \hline
    \end{tabular}
    \source{Elaborado pelo autor (2026)}
\end{table}

Com esses dados em mente, podemos compreender alguns elementos fundamentais:
poucos sistemas analisados são \textit{open-source}, o que indica que, em sua
grande maioria, são sistemas comerciais. A maior parte dos sistemas também
oferecem algum aspecto de pagamento, bloqueando funcionalidades importantes
atrás de algum tipo de modelo de subscrição. Existe também o fato de
personalização e de exportação de dados, que poucos ou nenhum sistema oferece,
o que faz com que o sistema não tenha portabilidade dos dados inseridos,
criando assim, um \textit{walled garden}, ou um jardim fechado -- uma técnica de
fazer com que o usuário não consiga migrar entre sistemas, prendendo-o a usar
aquela aplicação já que seus dados estão lá.

\subsection{Relação entre os dados coletados e o desenvolvimento}

A análise dos sistemas similares, quando considerada em conjunto com os dados
coletados por meio do questionário, revela lacunas e oportunidades relevantes
para o desenvolvimento de uma nova solução. Observa-se, por exemplo, que a
maior parte das ferramentas disponíveis no mercado segue modelos comerciais,
frequentemente com funcionalidades essenciais bloqueadas por sistemas de
pagamento ou assinaturas. Além disso, é comum que tais sistemas sejam fechados,
sem disponibilização de código-fonte ou possibilidade de exportação dos dados
inseridos, o que configura uma prática de \textit{walled garden}, dificultando
a migração entre plataformas e restringindo a autonomia dos usuários sobre suas
próprias informações.

Esses aspectos, quando contrastados com o interesse expressivo dos respondentes
por visualizar padrões emocionais ao longo do tempo e por maior personalização,
sugerem direções claras para o desenvolvimento de um sistema que não apenas
atenda a essas demandas, mas também adote princípios de transparência, abertura
e controle por parte do usuário. A seguir, serão discutidas as decisões de
projeto que nortearam a construção da ferramenta proposta, baseadas tanto nas
limitações identificadas nos sistemas existentes quanto nas necessidades e
expectativas manifestadas pelo público-alvo da pesquisa.

\chapter{Nexo}

O aplicativo proposto por esse trabalho, denominado "Nexo", se propõe a ser uma
aplicação para dispositivos móveis, de forma a preencher essas lacunas
demonstradas tanto pelo questionário, quanto também tentando preencher as
questões analisadas em outros sistemas similares. Como mencionado
anteriormente, o sistema é direcionado para pessoas com TAB, embora, ele também
possa atuar isoladamente como um aplicativo de registro de humor, e portanto, o
nome reflete essa escolha. A lógica é, de um aplicativo que consegue conectar
diversos pontos, que em um primeiro momento parecem soltos: sono, rotina,
práticas, humor, gatilhos.

A partir das demandas identificadas na pesquisa e dos princípios estabelecidos
para este projeto, delineia-se um sistema de monitoramento de humor centrado no
usuário, com ênfase em personalização, privacidade e utilidade reflexiva. O
funcionamento do sistema parte da premissa de que o humor não é uma variável
universal, mas sim composto por elementos subjetivos que variam entre
indivíduos. Assim, fundamentalmente para cada entrada no diário, o usuário pode
indicar que componentes de humor ele percebe em vigência naquele momento, como
``raiva'', ``tristeza'', ``felicidade'', etc., elementos que ele considera mais
representativos de seu estado emocional, permitindo um mapeamento mais fiel da
sua experiência.

O monitoramento, pode ser gerado em diversos momentos do dia, o que torna o
aplicativo mais flexível, ao invés de uma única entrada no dia. O usuário pode
também anotar para cada momento, algo relevante, através de escrita livre, e
usar uma escala que possa indicar o que ele percebe como indicadores dos níveis
de energia, estresse e ansiedade. Associado a isso, o sistema permitirá a
entrada de gatilhos e eventos significativos, sejam positivos ou negativos, que
poderão ser posteriormente correlacionados com os estados de humor, permitindo
identificar padrões ou ciclos. Esses dados podem ser correlacionados/vinculados
com ``práticas'' ou ``intervenções'', como fármacos, psicoterapia, atividades
físicas.

Outro ponto essencial é a exportação de dados em formatos acessíveis, como o
formato CSV \cite{Shafranovich_2005}, garantindo que o usuário tenha controle e
mobilidade sobre suas informações, algo frequentemente negligenciado em
sistemas fechados. Além disso, a ferramenta é por inteiro \textit{software}
livre, garantindo transparência de funcionamento, possibilidade de auditoria e
fomento a melhorias comunitárias, especialmente importantes no contexto de
tecnologias sensíveis à saúde mental. Essas características formam a base de um
sistema comprometido com o bem-estar, o respeito à individualidade e a
soberania dos dados pessoais.

Com base nas funcionalidades listadas, pode-se definir os requisitos funcionais
e não funcionais no qual o sistema deve necessitar.

\section{Classificação dos Requisitos}

Como observado por \cite[p.~58]{sommerville:2010}, ``Os requisitos de um sistema são
as descrições dos serviços que o sistema deve prestar e as restrições à sua
operação''. Portanto, requisitos são declarações articuladas descrevendo de
forma clara o que o sistema deve ser capaz de fazer para satisfazer as
necessidades postas anteriormente. São definidas em duas categorias: requisitos
funcionais que descrevem o comportamento e funções do sistema e requisitos não
funcionais que descrevem os critérios específicos que podem ser usados para
avaliar o bom funcionamento do sistema. A separação entre requisitos funcionais
e não funcionais permite uma estrutura clara para priorização, testes e
validação do projeto, respeitando tanto os objetivos técnicos quanto éticos da
proposta. Logo a seguir, serão abordados com mais detalhes cada um deles.

\subsection{Requisitos funcionais}

Os requisitos funcionais dizem respeito às funcionalidades centrais do sistema,
ou seja, declarações de tudo o que o sistema deve e está proposto a fazer e
reagir. Sommerville reforça esse aspecto:

\begin{quote}

  Os requisitos funcionais de um sistema descrevem o que ele deve fazer. Eles
  dependem do tipo de software a ser desenvolvido, de quem são seus possíveis
  usuários e da abordagem geral adotada pela organização ao escrever os
  requisitos. Quando expressos como requisitos de usuário, os requisitos
  funcionais são normalmente descritos de forma abstrata, para serem
  compreendidos pelos usuários do sistema.

  \cite[p.~59]{sommerville:2010}
\end{quote}

Portanto, diante das funcionalidades previstas, que foram brevemente listadas,
no Quadro 2 é apresentado os requisitos que o sistema deve fornecer:

\newpage

% TODO: Add this to the list of tables
\begin{longtable}{| p{.20\textwidth} | p{.30\textwidth} | p{.40\textwidth} |}
    \caption{Requisitos funcionais}
    \label{tab:requisitos_funcionais} \\
        \hline
        N° & Requisito & Descrição \\
        \hline
        \endfirsthead
        \hline
        N° & Requisito & Descrição \\
        \hline
        \endhead
        \endfoot

        [RF01] &
        Login &
        O sistema deve ser capaz de realizar login com email e senha \\

        \hline

        [RF02] &
        Redefinir senha &

        O sistema permite que o usuário redefina a sua senha, mediante informar
        o email usado \\

        \hline

        [RF03] &
        Cadastro &
        O sistema deve permitir que o usuário crie uma conta \\

        \hline

        [RF04] &
        Desconectar &
        O sistema permite que o usuário faça logout do sistema \\

        \hline

        [RF05] &
        Manter humor &
        O usuário deve poder registrar e atualizar seu estado emocional
        conforme a necessidade \\

        \hline

        [RF06] &
        Registrar gatilho &
        O sistema deve permitir associar sentimentos ou eventos específicos que
        influenciam o humor \\

        \hline

        [RF07] &
        Visualizar histórico &
        O usuário deve ter acesso aos registros anteriores de humor e demais
        tipos de dados como Gatilho, Sono, etc. \\

        \hline

        [RF08] &
        Visualizar detalhes do dia &
        O sistema deve permitir expandir e consultar dados específicos de um
        determinado dia \\

        \hline

        [RF09] &
        Visualizar registros de humor recente &
        O usuário deve poder visualizar os ultimos registros de humor como
        forma de ler \\

        \hline

        [RF10] &
        Manter registro de sono &
        O sistema deve permitir que o usuário possa cadastrar registros de sono \\

        \hline

        [RF11] &
        Manter ações de cuidado &
        O usuário deve poder cadastrar ações de cuidado (atividade, medicamentos, consultas) \\

        \hline

        [RF12] &
        Exportar dados &
        O sistema deve oferecer uma funcionalidade para exportar os dados
        registrados \\

        \hline

        [RF13] &
        Registrar que um medicamento foi tomado &
        O usuário pode marcar que um determinado medicamento foi tomado
        mediante que foi cadastrado \\

        \hline

        [RF14] &
        Associar humor à um gatilho &
        O sistema deve permitir que o usuário vincule um humor específico com
        um gatilho \\

        \hline

        [RF15] &
        Associar consulta com um humor ou gatilho &
        O sistema deve permitir vincular uma consulta específica a um
        gatilho ou humor registrado \\

        \hline

        [RF16] &
        Gerar um relatório com os últimos dados &
        O usuário deve selecionar um período e gerar/exportar um relatório com
        os últimos dados  \\

        \hline
        \multicolumn{3}{c}{\small Fonte: Elaborado pelo autor (2026)} \\
\end{longtable}

\subsection{Requirementos não funcionais}

De acordo com \cite[p.~89]{sommerville:2010}: ``Requisitos não funcionais são
restrições sobre os serviços ou funções oferecidas pelo sistema''. Ou seja,
definem os critérios de qualidade que o sistema deve seguir, abrangendo
aspectos como segurança, usabilidade, desempenho e manutenção. Neste projeto,
destacam-se especialmente a preservação da privacidade do usuário, a
transparência proporcionada pelo código aberto e a possibilidade de utilização
offline, o que garante acessibilidade mesmo em contextos de conectividade
limitada.

\begin{longtable}{| p{.20\textwidth} | p{.30\textwidth} | p{.40\textwidth} |}
    \caption{Requisitos não funcionais}
    \label{tab:requisitos_nao_funcionais} \\
        \hline
        N° & Requisito & Descrição \\
        \hline
        \endfirsthead
        \hline
        N° & Requisito & Descrição \\
        \hline
        \endhead
        \endfoot

        [RNF01] &
        Usabilidade &

        O sistema deve possuir uma interface intuitiva e acessível para
        usuários leigos em tecnologia, com navegação fluída e confirmações
        visuais após cada interação importante \\

        \hline

        [RNF02] &
        Desempenho &

        O tempo de resposta entre a entrada do usuário e a resposta do sistema
        deve ser inferior a 1 segundo; O carregamento do histórico semanal e
        visualizações de ciclos anteriores deve ocorrer em no máximo 2
        segundos. \\

        \hline

        [RNF03] &
        Segurança &

        Os dados do usuário devem ser armazenados de forma criptografada; O
        sistema deve exigir autenticação segura para acesso; A exportação de
        dados deve sempre exigir a confirmação do usuário \\
        \hline

        [RNF04] &
        Confiabilidade &

        O sistema deve garantir a persistência dos dados, mesmo em casos de
        falha de energia ou conexão; Os dados registrados devem ser armazenados
        automáticamente e sincronizados quando ouvir conexão \\

        \hline

        [RNF05] &
        Portabilidade &

        Todos os dados do usuário podem ser exportados a qualquer momento, em
        um formato de texto que permita que ele possa ler e interpretar,
        mediante compreensão mínima da linguagem/formato utilizado \\

        \hline

        [RNF06] &
        Privacidade &

        Nenhuma informação sensível do usuário deve ser compartilhada com
        terceiros e o sistema deve fornecer uma política de privacidade clara e
        acessível; Quaisquer alterações à mesma demanda que o usuário seja
        notificado \\

        \hline
        \multicolumn{3}{c}{\small Fonte: Elaborado pelo autor (2026)} \\
\end{longtable}

Com essas definições em mãos, pode-se avançar para a etapa de modelagem do
sistema, usando os requerimentos apresentados como fundação para o sistema.

\chapter{Modelagem}

De acordo com \cite{guedes:2022}, a modelagem de sistemas é uma etapa fundamental do
processo de desenvolvimento de software, pois permite representar graficamente
os principais elementos do sistema, seus comportamentos e interações. Essa
representação facilita a comunicação entre desenvolvedores, usuários e demais
envolvidos, servindo tanto como guia para implementação quanto como instrumento
de validação dos requisitos levantados. Através da modelagem, é possível
antever a estrutura do sistema, simular sua dinâmica e assegurar que ele
atenderá às necessidades específicas dos usuários.

Neste capítulo, são apresentados três tipos principais de diagramas da
\textit{Unified Modeling Language} (UML) empregados na modelagem do sistema
proposto: o diagrama de casos de uso, que descreve as funcionalidades esperadas
a partir da perspectiva do usuário; o diagrama de classes, que estrutura os
elementos do domínio em termos de atributos e relacionamentos; e o diagrama de
sequência, que detalha a comunicação entre os objetos ao longo do tempo. Esses
diagramas complementam-se mutuamente, proporcionando uma visão abrangente do
comportamento e da arquitetura do sistema.

\section{Casos de Uso}

O diagrama de casos de uso é o mais geral, informal e simples de compreender
\cite{guedes:2022}. A base dele é o levantamento de requisitos e a pesquisa
realizada, e serve como o pontapé para todo o processo de modelagem.

O diagrama de casos de uso, apresenta de forma simples, os atores, que são
representações das entidades que interagem com o sistema (sejam eles usuários
finais ou outros sistemas), e os casos de uso, que são as funcionalidades que o
sistema oferece. Ambos são conectados por linhas de associação. O sistema
também é delimitado por um retângulo, que contém os casos de uso
\cite{ibm:2021}.

Ao final desse processo, essa modelagem é transformada em documentação, e no
apêndice B deste documento, contém a mesma, oferecendo uma visão mais detalhada
para cada caso de uso, seus caminhos alternativos e caminhos de exceções.
Pode-se verificar que esse acaba por ser o diagrama mais basilar do projeto
inteiro \cite{guedes:2022}.

\begin{figure}[h]
\centering
\caption{Diagrama de Casos de Uso}
\label{fig:use-cases}
\includegraphics[width=0.9\linewidth]{casos-de-uso.png}
\source{O autor}
\end{figure}

Na Figura~\ref{fig:use-cases}, pode-se observar que apenas um ator interage
com o sistema, o próprio usuário final. Também visualiza-se que cada elipse,
se resume a um requisito funcional, e os casos de uso estão interligados por
setas nomeadas com a palavra \textit{''<<extend>>''} que indica conexões onde a
execução de um caso de uso adjacente pode ocorrer dado uma determinada
condição em adição ao caso de uso inicial \cite{ibm:2021}.

Permeando todas as interações no sistema, o caso de uso mais central, o
''Login'', cujo qual estabelece uma base de autenticação, de forma que garante
que apenas usuários autenticados possam interagir com o sistema.

No caso de uso ''Manter Humor'' representa a funcionalidade central do sistema,
permitindo que o usuário registre seu estado emocional ao longo do tempo. Essa
ação ocorre por meio da inserção direta de uma informação de humor, cujo qual o
usuário pode selecionar entre cinco humores ''basais'', além de elementos que
ajudem a compreender, como escalas indicativas dos respectivos níveis de
estresse, ansiedade e quais componentes do humor o usuário identifica com mais
presença. A manutenção do humor no sistema contempla tanto a criação quanto a
edição e a exclusão desses registros.

O caso de uso ''Registrar Intervenção'' contempla o cadastro de estratégias,
práticas ou sugestões que o usuário utiliza ou deseja utilizar para lidar com
oscilações emocionais. As ações de cuidados, podem incluir desde ações
práticas, como caminhar, até técnicas cognitivas ou expressivas, como
atividades criativas, medicamentos utilizados ou atendimentos com profissionais
da área. Essa funcionalidade contribui para que o sistema sirva como uma
central e colabora para construir um registro mais criterioso, visto que o
usuário reconstrói o que aconteceu naquele determinado ponto no tempo.

O caso de uso ''Registrar Gatilho'' visa permitir que o usuário registre
eventos, sejam eles pensamentos, situações, ou estímulos que tenham
influenciado significativamente seu estado emocional. O objetivo é construir um
histórico de causas associadas a variações no humor, de modo a fomentar a
identificação de padrões recorrentes ou específicos. Essa funcionalidade é
especialmente útil para usuários em contextos terapêuticos ou de
autoconhecimento, favorecendo a análise de relações entre circunstâncias
externas e reações internas.

O caso de ''Registro de sono'', permite o usuário, de forma simples, criar no
sistema um registro da quantidade de horas de sono, e adicionar alguma anotação
breve. Esse elemento corrobora com a construção de indicadores, tanto da média
de sono, quanto também de estimular uma correlação entre sono e os níveis de
energia informado em cada humor.

Por fim, o caso de uso ''Vincular humor a um gatilho'' permite ao usuário
estabelecer vínculos entre os gatilhos identificados e os registros de humor,
formando um mapeamento de ação e reação, ou causa e reação. Essa associação
pode servir tanto para consulta futura quanto para, em um futuro, gerar
indicadores. Ao conectar eventos e suas respostas, a funcionalidade fortalece o
papel do sistema como uma ferramenta alinhada com o EMA, gerando assim uma
narrativa emocional diária.

\section{Diagrama de Classes}

O diagrama de classes é um dos principais artefatos da UML, responsável por
evidenciar a estrutura estática do sistema, seus atributos, operações e
relacionamentos. Conforme \cite[p. 101]{guedes:2022}: ''O principal enfoque é
permitir a visualização das classes que o sistema e seus respectivos métodos e
atributos''. Dessa forma, esse diagrama constitui a base da modelagem
estrutural que posteriormente será implementada em código.

Sobre o diagrama em si, cada classe possui três partes: título, os atributos,
que representam as características da classe em si, e seus métodos, que são as
operações que o sistema pode operar nos objetos das classes \cite{guedes:2022}. Na
Figura 2, pode-se observar os elementos descritos acima:

\begin{figure}[h]
\centering
\caption{Diagrama de Classes}
\label{fig:classes}
\includegraphics[width=0.9\linewidth]{diagrama-classes.png}
\source{O autor}
\end{figure}

Essa representação apresenta alguns vários elementos importantes. Primeiro, um
acoplamento entre \textit{Usuário} e as demais classes. Isso é natural, e
desejado, considerando que é o único ator. Existe também a separação entre as
entidades como Gatilho, Humor, muito embora exista uma relação de associação
entre elas, através da classe denominada ''VinculoGatilhoHumor'', visto que ela
seria representada por uma tabela auxiliar ligando as duas, aliadas do atributo
de estimativa de impacto.

Necessário destacar aqui que a classe ''AcaoCuidado'' é uma classe abstrata,
representando uma generalização. Essa classe serve como um guarda-chuva para
representar diversos tipos de ações que o usuário pode realizar. O papel dela
aqui se torna guardar elementos em comum entre esses tipos de ações, seus
respectivos vínculos com gatilhos ou humores. A partir dela então, traça-se
outras sub-classes que contem um melhor contexto sobre a atividade em si, por
exemplo, na classe ''Consulta'' guarda-se dados como a duração da consulta, que
tipo de consulta, e potenciais anotações ou comentários em cima da mesma.
Entretanto, uma classe ''RegistroMedicamento'' não precisa guardar mais do que
a relação com qual medicamento e quando ele foi de fato tomado.

Por fim, existe diversos tipos de enumeradores que delimitam exatamente quais
tipos de valores podem ser escolhidos. Esses enumeradores servem também como um
mecanismo de validação, sendo movidos tanto para o próprio banco de dados e
posteriormente sendo implementados como forma de controle em outros pontos como
na API.

Ressalta-se, contudo, que o diagrama não contempla, padrões arquiteturais
complementares, como \textit{Service Layer}, \textit{Repository} e \textit{Data
Access Object (DAO)}, os quais são incorporados na etapa de implementação para
garantir uma separação adequada de responsabilidades e facilitar a manutenção e
escalabilidade do código, conforme as propostas de \citet{Fowler_2013} e
\citet{Gamma:2011}. Esses elementos são importantes, considerando os princípios
de legibilidade, simplicidade e responsabilidades, muitos organizados por
\citet{Martin_2012}, asseguram que o código produzido seja limpo e sustentável.

\input{chapters/design}
\chapter{Tecnologias}

Este capítulo apresenta as principais tecnologias empregadas no desenvolvimento
do aplicativo, organizadas de acordo com a camada da arquitetura em que
atuam: front-end, back-end e infraestrutura/persistência de dados.

\section{Front-end}

\subsection{React Native}

React Native é um \textit{framework} de código aberto, mantido pela Meta, que
permite o desenvolvimento de aplicativos móveis multiplataforma a partir de
uma única base de código escrita em JavaScript/TypeScript, utilizando
componentes que são traduzidos para elementos nativos tanto no iOS quanto no
Android \citep{reactnative:2026}.

A escolha do React Native se justifica pelo escopo reduzido da equipe de
desenvolvimento (um único desenvolvedor) e facilidade de conhecimento e
manutenção, sem trazer, para este projeto, benefícios que justifiquem esse
custo adicional. Além disso, o ecossistema maduro do React Native oferece
ampla disponibilidade de bibliotecas para funcionalidades já previstas no
aplicativo, como notificações locais e gráficos.

\subsection{Zustand}

Zustand é uma biblioteca leve de gerenciamento de estado para aplicações
React e React Native, que dispensa a estrutura verbosa de \textit{providers}
e \textit{reducers} tipicamente associada a bibliotecas como Redux
\citep{zustand:2026}.

Ela foi escolhida para gerenciar o estado efêmero da aplicação (por exemplo,
dados de formulários em preenchimento ou estados de navegação temporários)
por sua API minimalista e baixo \textit{overhead} de configuração, adequada a um
projeto de escopo individual em que a complexidade adicional de soluções
mais robustas de gerenciamento de estado não se justifica.

\section{Back-end}

\subsection{Express.js}

Express.js é um \textit{framework} minimalista para aplicações web em
Node.js, amplamente utilizado para a construção de APIs \textit{REST},
oferecendo uma camada fina de abstração sobre o roteamento de requisições HTTP e
\textit{middlewares} \citep{express:2026}.

Sua adoção se deve à simplicidade e à flexibilidade não opinativa do
\textit{framework}, que permite estruturar a API de acordo com as
necessidades específicas do projeto, sem impor convenções rígidas de
arquitetura, além de sua ampla documentação e maturidade no ecossistema
Node.js.

\subsection{Zod}

Zod é uma biblioteca de validação e definição de esquemas de dados para
TypeScript, que permite declarar formatos esperados de entrada (como corpos
de requisição) e inferir automaticamente os tipos estáticos correspondentes
\citep{zod:2026}.

Zod foi escolhido para validar os dados recebidos pela API antes de seu
processamento, garantindo que registros de humor, sono e demais entradas do
usuário estejam em conformidade com o formato esperado. A inferência
automática de tipos a partir dos esquemas evita a duplicação de definições de
tipo entre a camada de validação e o restante do código TypeScript do
back-end.

\subsection{Valkey e BullMQ}

Valkey é um banco de dados em memória, do tipo chave-valor, criado como
\textit{fork} de código aberto do Redis, mantido sob a governança da Linux
Foundation \citep{valkey:2026}. BullMQ é uma biblioteca para gerenciamento de
filas de tarefas em Node.js, construída sobre bancos de dados compatíveis com
o protocolo Redis \citep{bullmq:2026}.

A combinação das duas ferramentas foi adotada para o processamento paralelo
e assíncrono de tarefas que não precisam bloquear a resposta imediata da API,
como o disparo de notificações programadas e o processamento de indicadores
a partir do histórico de registros do usuário. A escolha do Valkey em
específico, e não do Redis diretamente, se deve à mudança de licenciamento do
Redis em 2024, que passou a restringir certos usos comerciais; o Valkey
mantém compatibilidade total de protocolo, permanecendo sob licença
permissiva de código aberto.

\section{Infraestrutura e Persistência de Dados}

\subsection{PostgreSQL}

PostgreSQL é um sistema gerenciador de banco de dados relacional, de código
aberto, reconhecido por sua conformidade com o padrão SQL, robustez e suporte
a tipos de dados avançados \citep{postgresql:2026}.

Sua escolha se justifica pela natureza estruturada e relacional dos dados do
domínio da aplicação (registros de humor associados a usuários, horários,
gatilhos e categorias), que se beneficia das garantias de integridade
referencial e das transações ACID oferecidas por um banco relacional
maduro, em detrimento de soluções não relacionais.

\subsection{Prisma}

Prisma é um \textit{Object-Relational Mapper} (ORM) para o ecossistema
Node.js/TypeScript, que gera automaticamente um cliente de acesso ao banco de
dados tipado a partir de um esquema declarativo, além de oferecer um sistema
de migrações versionadas \citep{prisma:2026}.

O Prisma foi escolhido por permitir a manutenção do esquema do banco de dados
de forma versionada e sincronizada com o código da aplicação, reduzindo a
chance de divergência entre a estrutura do banco e as entidades manipuladas
pelo back-end, além de prover checagem de tipos em tempo de compilação para
as consultas realizadas.

\section{Implantação e Automação}

\subsection{Docker e Docker Compose}

Docker é uma plataforma de conteinerização que permite empacotar uma
aplicação e suas dependências em uma unidade isolada e portátil, capaz de ser
executada de forma consistente em diferentes ambientes \citep{docker:2026}. O
Docker Compose é uma ferramenta complementar que permite definir e orquestrar
múltiplos contêineres -- e as relações entre eles -- por meio de um único
arquivo de configuração declarativo \citep{dockercompose:2026}.


No projeto, o Docker Compose é utilizado tanto no ambiente de desenvolvimento
quanto no ambiente de produção para orquestrar os serviços de banco de dados
(PostgreSQL), armazenamento em memória (Valkey) e da própria API, garantindo
que as versões e configurações desses serviços sejam idênticas entre as
máquinas de desenvolvimento e o servidor de produção.

O uso de \textit{healthchecks} em conjunto com a diretiva
\texttt{depends\_on} e a condição \texttt{service\_healthy} assegura
que a API só seja iniciada após o banco de dados e o Valkey estarem de
fato prontos para aceitar conexões, evitando falhas de inicialização
por condição de corrida (\textit{race condition}), como ilustrado no
Código~\ref{lst:healthcheck-compose}.

\begin{lstlisting}[caption={Configuração dUndefined referencee \textit{healthcheck} do Valkey e de dependência condicional da API no \texttt{docker-compose.yml}}, label={lst:healthcheck-compose}]
valkey:
  image: valkey/valkey:9.0-alpine
  healthcheck:
    test: ["CMD", "valkey-cli", "ping"]
    interval: 10s
    timeout: 5s
    retries: 5
    start_period: 5s

api:
  build: ./api
  depends_on:
    valkey:
      condition: service_healthy
    db:
      condition: service_healthy
\end{lstlisting}

\sourcecode{Elaborado pelo autor (2026)}

A adoção do Docker se justifica, no contexto de uma aplicação de saúde
mental de código aberto, pelo objetivo de garantir a reprodutibilidade do
ambiente de execução: qualquer pessoa que clone o repositório --- seja para
contribuir com o projeto, auditar seu funcionamento ou realizar sua própria
implantação --- consegue reproduzir o ambiente completo sem depender de
configurações manuais específicas do sistema operacional hospedeiro, o que
reforça o caráter \textit{self-service} do projeto.

\subsection{GitHub Actions}

GitHub Actions é uma plataforma de integração e entrega contínuas (CI/CD)
nativa do GitHub, que permite a definição de fluxos de trabalho automatizados
--- disparados por eventos como \textit{push} ou \textit{pull request} ---
descritos de forma declarativa em arquivos YAML \citep{githubactions:2026}.

No projeto, o fluxo de trabalho de CI/CD é dividido em três etapas
sequenciais e dependentes entre si. A primeira etapa (\texttt{test}) sobe um
serviço efêmero de PostgreSQL e executa a suíte de testes automatizados da
API a cada \textit{push} ou \textit{pull request} que modifique arquivos do
diretório da API.

A segunda etapa (\texttt{build-and-push}), condicionada ao
sucesso da primeira e restrita a eventos de \textit{push}, constrói a imagem
Docker de produção da API e a publica no GitHub Container Registry (GHCR),
com marcação automática pelo SHA do \textit{commit} e, condicionalmente, pela
tag \texttt{latest} quando o \textit{push} ocorre na \textit{branch}
\texttt{master}.

Por último, a terceira etapa (\texttt{deploy-production}), restrita à
\textit{branch} \texttt{master}, conecta-se ao servidor de produção via SSH e
realiza a atualização do serviço da API por meio da atualização da tag da
imagem no arquivo \texttt{.env} e da recriação do contêiner correspondente
via Docker Compose.

Esse processo elimina a necessidade de intervenção manual para a maior
parte do ciclo de implantação, restringindo a ação humana direta sobre o
servidor de produção a situações excepcionais, o que reduz a superfície de
erro humano e aumenta a confiabilidade do processo de entrega de novas
versões da aplicação.

\subsection{Ansible}

Ansible é uma ferramenta de automação de infraestrutura que permite
descrever, de forma declarativa e idempotente, o estado desejado de um ou
mais servidores, dispensando a necessidade de um agente instalado na máquina
gerenciada, uma vez que se comunica com ela via SSH \citep{ansible:2026}.

No escopo deste projeto, o papel do Ansible é ortogonal ao do Docker Compose
e ao do GitHub Actions: enquanto estes automatizam, respectivamente, a
orquestração dos serviços da aplicação e o ciclo de integração e entrega
contínuas, o Ansible é responsável exclusivamente pelo provisionamento
inicial do servidor de produção, executado uma única vez --- ou sempre que um
novo servidor precisar ser provisionado. Esse provisionamento consiste em:
(1) criar o usuário dedicado à aplicação; (2) adicionar a chave pública SSH
utilizada pelo GitHub Actions para se autenticar no servidor durante a etapa
de implantação; (3) criar os diretórios da aplicação com as permissões
adequadas; e (4) instalar o Docker no servidor. A cópia do arquivo de
variáveis de ambiente de produção (\texttt{.env.production}), por conter
segredos e credenciais sensíveis, é deliberadamente mantida como uma etapa
manual, não automatizada, e não é versionada no repositório.

A separação entre o provisionamento inicial (\textit{Ansible}), a orquestração de
execução (\textit{Docker Compose}) e a automação do ciclo de
entrega (\textit{GitHub Actions}) reflete uma decisão arquitetural:
cada ferramenta atua em uma camada distinta e bem delimitada do processo
de implantação, o que favorece a reprodutibilidade de todo o processo,
do zero, por qualquer pessoa que deseje hospedar sua própria instância da aplicação.

Essa preocupação é particularmente relevante no contexto de um
aplicativo de saúde mental de código aberto: ao reduzir a dependência
de serviços proprietários de implantação (\textit{walled gardens}) e
ao documentar e automatizar o processo de provisionamento, o projeto
favorece um modelo mais alinhado aos princípios de privacidade e de
autonomia sobre os dados dos usuários, permitindo que instituições ou
indivíduos preocupados com a confidencialidade de dados sensíveis de
saúde mental hospedem e controlem sua própria instância da aplicação.

\chapter{Desenvolvimento}

Este capítulo detalha o funcionamento de algumas funcionalidades
desenvolvidas durante o projeto. Apenas algumas funcionalidades mais
chaves são abordadas neste capítulo, como o registro de humor, vinculo
de um gatilho, indicadores, e o registro de atividades.

\section{Processo de Login}

Logo ao abrir o aplicativo, o usuário é redirecionado para a tela de
login. A tela em si é relativamente simples, e contém os campos de
email e senha. Existe aqui um elemento proposital: uma forma de
garantir um certo nível de anonimato, em que o usuário não precisa se
identificar com um nome ou número de telefone, tampouco ativar a conta
através de uma plataforma de terceiro, como Google ou Facebook.

Isso se dá, em especial, por dois motivos. Primeiro, o uso de OAuth
exigiria uma conexão com plataformas de terceiro, que poderiam então
rastrear quais usuários têm acesso a qual aplicativo — algo que é
fartamente documentado na literatura sobre os mecanismos usados por
tais plataformas para rastrear e monitorar o acesso a dados de
usuários, com o fim de vender e monetizar essas informações.

Essa preocupação não é hipotética: um estudo que interceptou o tráfego de rede de
36 dos aplicativos mais populares de saúde mental (depressão e cessação de
tabagismo) encontrou transmissão de dados a pelo menos um terceiro em 92\% deles,
sendo que 81\% enviavam dados especificamente a serviços operados por Google ou
Facebook — e apenas 59\% desses casos declaravam esse compartilhamento em política
de privacidade \cite{Huckvale2019}.

Um estudo posterior, focado exclusivamente em 27 aplicativos de saúde
mental mais populares da Google Play Store, confirmou o mesmo padrão e
o caracterizou formalmente como uma ameaça de \textit{linkability}
(vinculabilidade) e \textit{identifiability} (identificabilidade): mesmo quando os
identificadores enviados a terceiros são pseudônimos (como \textit{advertising IDs}
ou hashes), eles ainda permitem que provedores de publicidade associem o
comportamento de um usuário através de diferentes aplicativos e serviços
\cite{Iwaya:22}. Como os próprios autores resumem, a mera vinculação de um
usuário a um aplicativo de saúde mental já pode revelar que essa pessoa possui
uma condição psicológica — seja ansiedade, depressão, ou, no caso de um
aplicativo como o Nexo, Transtorno Afetivo Bipolar —, tornando usuários desses
aplicativos particularmente vulneráveis a esse tipo de exposição
\cite{Iwaya:22}.

No caso do Nexo, essa cadeia de inferência é especialmente
sensível: a presença do aplicativo na lista de apps instalados de um usuário,
associada ao seu identificador publicitário, permitiria a uma plataforma como a
Meta inferir uma possível condição de saúde mental do usuário para fins de
segmentação de anúncios, sem que este tenha conhecimento ou consentimento
explícito disso.

Segundo, e retomando a proposta de oferecer ao usuário a possibilidade
de \textit{self-service}, aquele que decidisse manter uma instância
própria do aplicativo não precisaria de uma conta em uma plataforma de
terceiro, nem se cadastrar nelas.

\begin{figure}[H]
\centering
\caption{Tela de login}
\label{fig:screenshot-login}
\includegraphics[width=0.70\textwidth,height=0.60\textheight,keepaspectratio]{screenshot-login.jpeg}
\source{O autor, 2026}
\end{figure}

Por fim, destaca-se aqui que, conforme mencionado anteriormente na
modelagem de casos de uso (figura \ref{fig:use-cases}), ao realizar o
login pela primeira vez é enviado um email com um código de
verificação, seguindo a lógica dos \textit{magic links}. Isso adiciona
uma camada extra de segurança, garantindo que apenas o usuário que
solicitou o login possa acessar a conta.

\section{Página de Consulta de Registros}

A tela inicial é a tela nomeada apenas como ``Novo''. A proposta aqui
é ser a tela de cadastro rápido de um humor, e também é a tela que
visualmente mais destoa das demais. O usuário é apresentado com três
seções distintas:

\begin{itemize}
    \item Cadastro rápido de um humor
    \item Entre ontem e hoje
    \item Registros recentes
\end{itemize}

Respectivamente: uma seção em que o usuário é convidado a ``escolher
um humor'' e iniciar o processo de registro de uma nova entrada; a
seção em que são apresentados os indicadores de nível de energia
acumulado entre o dia anterior e o dia atual, além da média de sono do
dia comparada com a do dia anterior; e, por fim, a listagem dos dez
últimos registros de humor cadastrados pelo usuário.

\begin{figure}[H]
\centering
\caption{Tela de Consulta de Registros}
\label{fig:screenshot-consulta-registros}
\includegraphics[width=0.70\textwidth,height=0.60\textheight,keepaspectratio]{screenshot-initial-page.jpeg}
\source{O autor, 2026}
\end{figure}

Aqui já é introduzido ao usuário um padrão de organização e hierarquia:
todos os elementos centrais são organizados em \textit{cards}, exibidos
em no máximo duas colunas. Também é introduzido, caso o usuário já
possua registros de humor cadastrados, o padrão pelo qual esses
registros são resumidos em todas as seções do aplicativo: geralmente
um ícone, um título e alguma informação adicional em destaque.

\section{Criação de Registros de Humor}

Nesta seção, o usuário pode criar novos registros de humor, escolhendo
entre uma lista de opções pré-definidas. O humor selecionado na tela
anterior é mostrado aqui, e o usuário é convidado a julgar, em uma
escala de 1 a 10, como se sente em relação aos níveis de energia,
estresse e ansiedade, respectivamente.

\begin{figure}[H]
\centering
\caption{Criação de Registros de Humor}
\label{fig:screenshot-consulta-humor}
\includegraphics[width=0.70\textwidth,height=0.60\textheight,keepaspectratio]{screenshot-create-mood.jpeg}
\source{O autor, 2026}
\end{figure}

Naturalmente, é uma questão difícil de responder — e, como o próprio
aplicativo procura deixar claro, é importante que seja um momento
reflexivo, não necessariamente preciso.

O usuário pode adicionar uma anotação sobre aquele registro em
específico, padrão que se repete em outras telas do aplicativo.

Por fim, o usuário pode escolher o que compõe aquele registro, ou
seja, selecionar emoções e estados/sentimentos que sente naquele
momento. É apresentada uma listagem mista desses elementos, na qual o
usuário seleciona os que mais se aplicam a ele e atribui um nível de
intensidade para cada um. Aqui a escala é menor e mais simples, visando
facilitar a escolha e o julgamento do usuário, entre ``Leve'',
``Moderado'' e ``Intenso'', conforme a figura
\ref{fig:screenshot-add-components}.

\begin{figure}[H]
\centering
\caption{Adicionando sentimentos/emoções}
\label{fig:screenshot-add-components}
\includegraphics[width=0.70\textwidth,height=0.60\textheight,keepaspectratio]{screenshot-add-components.jpeg}
\source{O autor, 2026}
\end{figure}

Ao concluir esse processo e salvar, é apresentada ao usuário uma tela
de confirmação, na qual ele pode, opcionalmente, vincular o registro a
um gatilho específico — que pode ter ocorrido antes de ele salvar o
registro — ou criar um novo gatilho. O projeto novamente se ancora no
EMA, enriquecendo o registro de humor com mais contexto sobre o que
aconteceu.

\begin{figure}[H]
\centering
\caption{Vinculando um humor com um gatilho}
\label{fig:screenshot-link-mood-trigger}
\includegraphics[width=0.70\textwidth,height=0.60\textheight,keepaspectratio]{screenshot-link-mood-with-trigger.jpeg}
\source{O autor, 2026}
\end{figure}

Em seguida, o usuário pode indicar o impacto que aquele gatilho
exerceu sobre o registro de humor em questão. Novamente, trata-se de
uma estimativa subjetiva, sem pretensão de precisão.

\begin{figure}[H]
\centering
\caption{Avaliando o impacto de um gatilho no humor}
\label{fig:screenshot-assess-trigger-mood-impact}
\includegraphics[width=0.70\textwidth,height=0.60\textheight,keepaspectratio]{screenshot-link-mood-assess-impact.jpeg}
\source{O autor, 2026}
\end{figure}

Por fim, vale mencionar que, após esse processo, é enfileirado
internamente um \textit{background job} que recalcula, sob demanda, o
indicador de energia daquele dia, apresentado na figura
\ref{fig:screenshot-consulta-registros}. Diferentemente deste, os
indicadores semanais — como a correlação entre sono e energia,
descrita na seção \ref{sec:insights} — não são recalculados a cada
registro individual, sendo processados apenas pelo job semanal
agendado (seção \ref{sec:insights}).

\section{Indicadores}\label{sec:insights}

Nessa tela é apresentado um conjunto de indicadores construídos a partir
dos registros de humor, sono e gatilho do usuário. Diferente das demais
telas do aplicativo, que operam sobre dados em tempo real, os indicadores
são artefatos pré-computados: um \textit{worker} assíncrono (\textit{BullMQ})
percorre semanalmente todos os usuários ativos e enfileira, para cada
um, um conjunto de \textit{jobs} — um por tipo de indicador — que
processam a janela dos últimos sete dias e persistem o resultado em uma
tabela dedicada. Vale notar que nem todos os indicadores seguem esse
fluxo agendado: os indicadores diários (energia e sono) são recalculados
sob demanda, no momento em que o usuário registra um novo humor ou
sessão de sono, enquanto os indicadores semanais dependem
exclusivamente do \textit{job} de \textit{fan-out} de segunda-feira.

O aplicativo móvel apenas consulta esse resultado já
calculado, o que evita repetir cálculos custosos a cada abertura da tela
e desacopla a lógica analítica do restante da API.

\begin{figure}[H]
\centering
\caption{Tela de indicadores}
\label{fig:screenshot-insights}
\includegraphics[width=0.70\textwidth,height=0.60\textheight,keepaspectratio]{screenshot-insights.jpeg}
\source{O autor, 2026}
\end{figure}

% TODO: referenciar EMA / momentary assessment como justificativa da janela de 7 dias

São três os indicadores semanais exibidos: tendência de humor, correlação
entre sono e energia, e padrão de gatilhos.

\begin{itemize}
    \item \textbf{Tendência de humor}: compara a média de humor da
    primeira metade da semana com a da segunda metade, expressando a
    diferença como uma variação percentual. O resultado é classificado
    como melhora, piora ou estabilidade a partir de uma margem de $\pm10\%$.

    \item \textbf{Correlação entre sono e energia}: relaciona a duração
    do sono de uma noite com o nível de energia relatado na manhã
    seguinte, através do coeficiente de correlação de Pearson. Exige um
    mínimo de observações pareadas para ser calculado — caso contrário,
    o indicador simplesmente não é exibido.

    \item \textbf{Padrão de gatilhos}: identifica a categoria de gatilho
    mais frequente na semana e expressa sua predominância como percentual
    do total de registros.
\end{itemize}

É intencional, na figura \ref{fig:screenshot-insights}, que o bloco de
``Energia'' apareça vazio: o cálculo da correlação depende da existência
simultânea de registros de sono \emph{e} de humor em quantidade mínima,
e no período retratado o usuário ainda não havia acumulado dados
suficientes de ambos os tipos. Optou-se por não exibir uma correlação
calculada com poucos pontos, preferindo comunicar a ausência de dado a
apresentar um número pouco confiável.

O bloco maior no topo da tela, o ``Resumo Semanal'', condensa os
indicadores mais centrais — média de sono do período e um rótulo de
humor (``Estável'', no exemplo). Esse rótulo, no entanto, é o elemento
mais delicado de toda a funcionalidade. Embora exista uma fórmula capaz
de produzir uma classificação de estabilidade a partir da variância dos
registros de humor, atribuir esse tipo de afirmação a um usuário —
especialmente em um aplicativo que é utilizado por pessoas com
Transtorno Afetivo Bipolar, para quem uma leitura equivocada de
``estabilidade'' pode mascarar justamente o início de um episódio —
exige um embasamento mais cuidadoso do que uma regra estatística
simples. Por esse motivo, o rótulo apresentado nesta versão do
aplicativo é estático (não deriva de fato do cálculo de estabilidade) e
a funcionalidade pode ser inteiramente desabilitada pelo usuário nas
configurações. A decisão de manter ou não uma afirmação automatizada
desse tipo, e sob quais critérios, permanece um ponto em aberto que
mereceria revisão de literatura específica antes de ser tratado como
funcionalidade definitiva.

% TODO: literatura sobre riscos de auto-monitoramento automatizado de humor / mood-tracking apps e TAB

\section{Histórico}


Nessa tela é apresentado o histórico consolidado de registros de humor,
sono, gatilho e ações de cuidado (como sessões de terapia), organizados
cronologicamente e agrupados por dia, conforme a figura
\ref{fig:screenshot-history}. O usuário pode ainda filtrar a listagem por
categoria (``Humor'', ``Sono'', ``Gatilho'' e ``Ações de Cuidado''),
reduzindo o volume de informação exibida.

\begin{figure}[H]
\centering
\caption{Tela de Histórico}
\label{fig:screenshot-history}
\includegraphics[width=0.70\textwidth,height=0.60\textheight,keepaspectratio]{screenshot-history.jpeg}
\source{O autor, 2026}
\end{figure}

Como os quatro tipos de dado residem em recursos distintos da API, a
tela consulta os respectivos \textit{endpoints} em paralelo e, em
seguida, processa o resultado no próprio aplicativo. Para isso, é
utilizado um pequeno serviço que atua como \textit{mapper}, transformando
os dados de cada \textit{endpoint} em um formato comum, apresentado no
código \ref{lst:history-card}.

\newpage

\begin{lstlisting}[caption={Formato comum da tela de histórico}, label={lst:history-card}]
export type HistoryCard = {
  id: string; // "<category>-<id>" to avoid collisions
  category: HistoryCategory;
  subtype?: CareActionSubtype; // only set when category === "care_action"
  timestamp: string; // ISO string -- used for sorting and display
  title: string;
  summary: string;
  badge?: HistoryBadge;
  raw: MoodEntry | SleepRecord | Trigger | CareAction;
};
\end{lstlisting}

Após a normalização, os quatro conjuntos de dados são combinados em uma
única lista ordenada pelo campo \texttt{timestamp} e reagrupados por dia
para exibição, conforme visto na figura \ref{fig:screenshot-history}, nos
cabeçalhos ``terça-feira, 14 de julho'' e ``segunda-feira, 13 de julho''.
Essa unificação é proposital: ao invés de apresentar quatro telas
separadas por tipo de dado, o histórico busca reconstruir uma narrativa
única do dia do usuário, aproximando-se do princípio de \textit{Ecological
Momentary Assessment} já mencionado anteriormente — o contexto de um
registro de humor (por exemplo, o gatilho ``Interno'' registrado três
minutos antes, na mesma figura) fica visualmente próximo do evento que
possivelmente o motivou.

A propriedade \texttt{raw} armazena o dado original, permitindo que a
listagem redirecione o usuário para a tela de detalhe correta ao tocar em
um card, independentemente da categoria. Já a propriedade \texttt{badge}
reaproveita o mesmo mecanismo já descrito na seção de indicadores: para
um registro de sono, sinaliza se a duração foi insuficiente; para um
registro de humor, destaca a dimensão mais relevante daquele registro
(como ``Ansiedade alta'' na figura \ref{fig:screenshot-history}); para um
gatilho, indica sua categoria, e assim por diante.

\FloatBarrier


\bibliography{monografia}
\bibliographystyle{abnt}

%%%%%%%%%%%%%%%%%%%%%%%%%%%%%%%%%%%%%%%%%%%%%%%%%%%%%%%%%%%%

% NÃO MEXER - RELACIONADO À FORMATAÇÃO DE ANEXOS/APÊNDICES
\makeatletter%
\renewcommand*{\@seccntformat}[1]{\csname the#1\endcsname\hspace{0.2cm}}%
\makeatother%

%%%%%%%%%%%%%%%%%%%%%%%%%%%%%%%%%%%%%%%%%%%%%%%%%%%%%%%%%%%%%%

% Adiciona um apêndice
\append{Documentação de Casos de Uso}

\input{extra/usecases}

\append{Resultados da Pesquisa}

\input{extra/survey-results}

% Adiciona um anexo
% \annex{Verão}

\end{document}
